\section{Introducci\'on}

El presente proyecto consiste en una aplicaci\'on m\'ovil multiplataforma orientada
a la consulta, visualizaci\'on y an\'alisis de informaci\'on deportiva de LaLiga. El
sistema permite explorar temporadas, equipos, jugadores y partidos, consultar
estad\'isticas hist\'oricas, gestionar favoritos, ejecutar simulaciones y acceder a
herramientas de an\'alisis predictivo.

La arquitectura l\'ogica se divide en cuatro componentes:

\begin{itemize}
  \item Una aplicaci\'on m\'ovil desarrollada con React Native y Expo.
  \item Una API REST desarrollada con Node.js y Express.
  \item Una base de datos PostgreSQL.
  \item Un m\'odulo anal\'itico en Python encargado de ejecutar la simulaci\'on
        Monte Carlo.
\end{itemize}

Adem\'as, el sistema incorpora un asistente conversacional basado en Gemini 2.5
Flash. Este asistente transforma preguntas en lenguaje natural en consultas SQL
controladas, ejecuta dichas consultas sobre PostgreSQL y genera una respuesta
comprensible para el usuario.

Desde el punto de vista del despliegue, el sistema se distribuir\'a mediante tres
contenedores Docker:

\begin{itemize}
  \item \texttt{database}, encargado de ejecutar PostgreSQL e inicializar el
        esquema y los datos.
  \item \texttt{backend}, que contendr\'a Node.js, Express, Python y las
        dependencias necesarias para la simulaci\'on Monte Carlo.
  \item \texttt{frontend}, responsable de ejecutar Expo y el servidor de
        desarrollo Metro Bundler.
\end{itemize}

El m\'odulo Python se integrar\'a dentro del contenedor del backend porque Express
lo ejecuta directamente como un proceso s\'incrono. Debido a que se utiliza para
una \'unica operaci\'on anal\'itica, no se considera necesario desplegarlo como un
microservicio independiente.

El objetivo de este manual es estandarizar el entorno de desarrollo y despliegue
para facilitar la ejecuci\'on, el mantenimiento y la ampliaci\'on del sistema,
conservando la separaci\'on de responsabilidades entre la interfaz m\'ovil, la API,
la persistencia y los servicios anal\'iticos.

\section{Preparaci\'on del entorno de trabajo}

La infraestructura del proyecto se contenerizar\'a mediante Docker, evitando que
el desarrollador tenga que instalar localmente PostgreSQL, Node.js, Python o sus
dependencias. La arquitectura estar\'a formada por tres contenedores comunicados
mediante una red interna de Docker.

\begin{table}[H]
\centering
\caption{Distribuci\'on de los servicios en contenedores.}
\label{tab:contenedores}
\begin{tabular}{|p{2.6cm}|p{8.2cm}|p{3cm}|}
\hline
\textbf{Contenedor} & \textbf{Contenido} & \textbf{Acceso} \\
\hline
\texttt{database} & PostgreSQL y volcado definitivo de la base de datos & Solo red interna \\
\hline
\texttt{backend} & Node.js, Express, Python y simulador Monte Carlo & Puerto 3000 \\
\hline
\texttt{frontend} & Expo y Metro Bundler & Puerto 8081 \\
\hline
\end{tabular}
\end{table}

La aplicaci\'on Expo Go no se ejecuta dentro de Docker, sino en un dispositivo
f\'isico o emulador. Esta se conecta al servidor Metro para descargar la aplicaci\'on
y directamente al backend para realizar las peticiones HTTP. El recorrido de las
comunicaciones es el siguiente:

\begin{lstlisting}
Expo Go ------------------> Metro Bundler :8081
   |
   +----------------------> Express :3000
                                  |
                                  +--> PostgreSQL :5432
                                  |
                                  +--> Python Monte Carlo
\end{lstlisting}

PostgreSQL no expondr\'a su puerto a la m\'aquina anfitriona, salvo que resulte
necesario durante las tareas de administraci\'on. El backend acceder\'a a la base de
datos mediante el nombre del servicio Docker, \texttt{database}, y el puerto
interno 5432.

Por su parte, los puertos 3000 y 8081 s\'i deber\'an publicarse, ya que Expo Go
necesita acceder tanto al backend como al servidor Metro desde el exterior de la
red Docker.

Los requisitos principales de la m\'aquina anfitriona se recogen en la
Tabla~\ref{tab:requisitos-host}.

\begin{table}[H]
\centering
\caption{Requisitos de software del entorno de desarrollo.}
\label{tab:requisitos-host}
\begin{tabular}{|p{3.4cm}|p{5cm}|p{4.6cm}|}
\hline
\textbf{Componente} & \textbf{Herramienta} & \textbf{Versi\'on m\'inima} \\
\hline
Contenedores & Docker Engine & 24.0 o superior \\
\hline
Orquestaci\'on & Docker Compose v2 & 2.20 o superior \\
\hline
Control de versiones & Git & 2.x \\
\hline
Cliente m\'ovil & Expo Go & Compatible con SDK 54 \\
\hline
Emulaci\'on opcional & Android Studio & Compatible con Expo SDK 54 \\
\hline
\end{tabular}
\end{table}

Aunque estas herramientas no tendr\'an que instalarse en el equipo anfitri\'on,
las im\'agenes Docker utilizar\'an internamente los entornos indicados en la
Tabla~\ref{tab:entornos-internos}.

\begin{table}[H]
\centering
\caption{Entornos internos previstos para los contenedores.}
\label{tab:entornos-internos}
\begin{tabular}{|p{2.6cm}|p{10.7cm}|}
\hline
\textbf{Contenedor} & \textbf{Entorno interno} \\
\hline
\texttt{database} & PostgreSQL 18.1 \\
\hline
\texttt{backend} & Node.js 20.19.4 o superior y Python 3.10 o superior \\
\hline
\texttt{frontend} & Node.js 20.19.4 o superior, React Native 0.81.5 y Expo SDK 54 \\
\hline
\end{tabular}
\end{table}

El contenedor \texttt{backend} instalar\'a tanto las dependencias npm de Express
como las dependencias Python del simulador: NumPy, pandas y psycopg2. Para ello,
ser\'a necesario incorporar al proyecto un archivo \texttt{requirements.txt}.

Las variables de entorno se distribuir\'an entre los servicios de la siguiente
forma:

\begin{itemize}
  \item \texttt{frontend}: \texttt{EXPO\_PUBLIC\_API\_URL}.
  \item \texttt{backend}: \texttt{DB\_USER}, \texttt{DB\_HOST},
        \texttt{DB\_NAME}, \texttt{DB\_PASSWORD}, \texttt{DB\_PORT},
        \texttt{GEMINI\_API\_KEY} y \texttt{MONTECARLO\_PYTHON\_PATH}.
  \item \texttt{database}: nombre de la base de datos, usuario y contrase\~na de
        PostgreSQL.
\end{itemize}

Dentro de la red Docker, el backend utilizar\'a \texttt{database} como valor de
\texttt{DB\_HOST}. Sin embargo, \texttt{EXPO\_PUBLIC\_API\_URL} deber\'a contener
la direcci\'on IP de la m\'aquina anfitriona y el puerto publicado del backend, ya
que un dispositivo f\'isico no puede resolver los nombres internos de Docker. Por
ejemplo:

\begin{lstlisting}
EXPO_PUBLIC_API_URL=http://192.168.1.100:3000/api
\end{lstlisting}

Una vez creados los archivos de contenerizaci\'on, la infraestructura completa se
iniciar\'a desde la ra\'iz del repositorio mediante el siguiente comando:

\begin{lstlisting}
docker compose up --build
\end{lstlisting}

\section{Estructura del repositorio}

El proyecto utiliza una estructura de monorepositorio en la que se almacenan
conjuntamente la aplicaci\'on m\'ovil, la API REST, el acceso a la base de datos y
el script de simulaci\'on. Una vez incorporada la contenerizaci\'on, la estructura
prevista ser\'a la mostrada en la Figura~\ref{fig:estructura-repositorio}.

\begin{figure}[H]
  \centering
  \begin{minipage}{0.92\textwidth}
\begin{lstlisting}
TFG_LaLiga/
|
|-- backend/
|   |-- src/
|   |   |-- config/
|   |   |   |-- db.js
|   |   |   |-- gemini.js
|   |   |   `-- swagger.js
|   |   |-- controllers/
|   |   |-- prompts/
|   |   |-- routes/
|   |   |-- scripts/
|   |   |   `-- simulacion_montecarlo_laliga.py
|   |   |-- services/
|   |   |-- utils/
|   |   `-- app.js
|   |-- Dockerfile
|   |-- requirements.txt
|   `-- package.json
|
|-- frontend/
|   |-- src/
|   |   |-- assets/
|   |   |-- components/
|   |   |-- context/
|   |   |-- screens/
|   |   `-- theme/
|   |-- App.js
|   `-- Dockerfile
|
|-- database/
|   `-- init/
|       `-- dump.sql
|
|-- index.js
|-- app.json
|-- babel.config.js
|-- package.json
|-- docker-compose.yml
|-- .dockerignore
|-- .env.example
`-- .gitignore
\end{lstlisting}
  \end{minipage}
  \caption{Estructura prevista del repositorio contenerizado.}
  \label{fig:estructura-repositorio}
\end{figure}

El archivo \texttt{docker-compose.yml} definir\'a los tres servicios, la red
interna, los puertos publicados, las variables de entorno, los vol\'umenes y las
dependencias de arranque.

El directorio \texttt{database/init/} contendr\'a el volcado SQL definitivo. Este
archivo se ejecutar\'a autom\'aticamente durante la primera inicializaci\'on del
contenedor PostgreSQL.

El archivo \texttt{Dockerfile} del backend instalar\'a Node.js y Python en una
misma imagen. Tambi\'en copiar\'a el c\'odigo del servidor, instalar\'a las
dependencias npm y procesar\'a el archivo \texttt{requirements.txt}. De esta
forma, Express podr\'a invocar el script de Monte Carlo directamente sin recurrir
a otro servicio de red.

El archivo \texttt{Dockerfile} del frontend instalar\'a las dependencias de Expo
y ejecutar\'a Metro Bundler escuchando en el puerto 8081.

Dentro del backend, el directorio \texttt{config/} contendr\'a la conexi\'on con
PostgreSQL, la configuraci\'on del cliente de Gemini y la definici\'on de
Swagger/OpenAPI. La carpeta \texttt{controllers/} implementar\'a la l\'ogica
asociada a usuarios, temporadas, equipos, jugadores, partidos, favoritos,
b\'usqueda, chatbot y Data Mining. Las rutas de la API permanecer\'an agrupadas
bajo el prefijo \texttt{/api}.

En el frontend, \texttt{App.js} configurar\'a la navegaci\'on, restaurar\'a la
sesi\'on almacenada y registrar\'a los proveedores globales. Las vistas completas
se organizar\'an en \texttt{screens/}, mientras que \texttt{components/} reunir\'a
los elementos reutilizables y los componentes especializados por dominio
funcional.

\section{Stack tecnol\'ogico}

Las principales tecnolog\'ias utilizadas por el sistema se recogen en la
Tabla~\ref{tab:stack-tecnologico}.

\begin{table}[H]
\centering
\caption{Stack tecnol\'ogico del proyecto (I).}
\label{tab:stack-tecnologico}
\begin{tabular}{|p{2.8cm}|p{4cm}|p{6cm}|}
\hline
\textbf{Capa} & \textbf{Tecnolog\'ia} & \textbf{Prop\'osito} \\
\hline
Contenerizaci\'on & Docker Engine & Aislamiento y ejecuci\'on de los servicios. \\
\hline
Orquestaci\'on & Docker Compose & Configuraci\'on conjunta del frontend, el backend y la base de datos. \\
\hline
Cliente & React 19.1 & Construcci\'on declarativa de la interfaz. \\
\hline
Cliente & React Native 0.81.5 & Desarrollo de la aplicaci\'on m\'ovil multiplataforma. \\
\hline
Cliente & Expo SDK 54 & Entorno de desarrollo y ejecuci\'on mediante Expo Go. \\
\hline
Desarrollo m\'ovil & Metro Bundler & Empaquetado y distribuci\'on del c\'odigo JavaScript. \\
\hline
Navegaci\'on & React Navigation 7 & Navegaci\'on mediante pila, men\'u lateral y pesta\~nas. \\
\hline
Persistencia local & AsyncStorage 2.2 & Almacenamiento de la sesi\'on y las preferencias del usuario. \\
\hline
Estado global & React Context & Gesti\'on de favoritos y del tema visual. \\
\hline
Visualizaci\'on & Chart Kit, Gifted Charts, Victory Native y SVG & Representaci\'on de estad\'isticas, gr\'aficas y dashboards. \\
\hline
\end{tabular}
\end{table}

\begin{table}[H]
\centering
\caption{Stack tecnol\'ogico del proyecto (II).}
\label{tab:stack-tecnologico-2}
\begin{tabular}{|p{2.8cm}|p{4cm}|p{6cm}|}
\hline
\textbf{Capa} & \textbf{Tecnolog\'ia} & \textbf{Prop\'osito} \\
\hline
Servidor & Node.js & Entorno de ejecuci\'on del backend. \\
\hline
API & Express 5.2.1 & Implementaci\'on de la API REST. \\
\hline
Configuraci\'on & dotenv y CORS & Variables de entorno y comunicaci\'on entre cliente y servidor. \\
\hline
Documentaci\'on & Swagger y OpenAPI 3.0 & Documentaci\'on y prueba interactiva de los endpoints. \\
\hline
Persistencia & PostgreSQL 18.1 & Almacenamiento de datos deportivos, usuarios y resultados anal\'iticos. \\
\hline
Acceso a datos & node-postgres (\texttt{pg}) & Pool de conexiones y ejecuci\'on de consultas SQL desde Express. \\
\hline
Seguridad & bcrypt & Cifrado y comprobaci\'on de contrase\~nas. \\
\hline
Inteligencia artificial & Gemini 2.5 Flash & Generaci\'on controlada de SQL y respuestas en lenguaje natural. \\
\hline
Anal\'itica & Python 3.10, NumPy y pandas & Procesamiento num\'erico empleado por la simulaci\'on. \\
\hline
Simulaci\'on & M\'etodo de Monte Carlo & Estimaci\'on probabil\'istica de la clasificaci\'on final. \\
\hline
Integraci\'on Python & psycopg2 & Comunicaci\'on entre los procesos Python y PostgreSQL. \\
\hline
\end{tabular}
\end{table}

La contenerizaci\'on descrita representa la arquitectura final prevista. En el
estado actual del proyecto todav\'ia deben incorporarse los archivos
\texttt{Dockerfile}, \texttt{docker-compose.yml}, \texttt{.dockerignore} y
\texttt{requirements.txt}, as\'i como trasladar el volcado definitivo al
directorio de inicializaci\'on de PostgreSQL.

\section{Instrucciones para la construcci\'on y puesta en marcha}

Una vez implementada la contenerizaci\'on descrita en los apartados anteriores,
la construcci\'on y puesta en marcha del sistema se realizar\'a de forma conjunta
mediante Docker Compose. El desarrollador solamente necesitar\'a tener instalados
Docker, Docker Compose, Git y Expo Go, ya que Node.js, PostgreSQL, Python y sus
dependencias se ejecutar\'an dentro de los contenedores correspondientes.

El primer paso consiste en obtener el c\'odigo fuente y acceder al directorio ra\'iz
del proyecto:

\begin{lstlisting}
git clone https://github.com/rafaPG03/TFG_LaLiga.git
cd TFG_LaLiga
\end{lstlisting}

\subsection{Configuraci\'on de las variables de entorno}

Antes de construir las im\'agenes se deber\'a crear el archivo de configuraci\'on
a partir de la plantilla incluida en el repositorio. En Linux o macOS se utilizar\'a
el siguiente comando:

\begin{lstlisting}
cp .env.example .env
\end{lstlisting}

En Windows se podr\'a realizar la misma operaci\'on desde PowerShell:

\begin{lstlisting}
Copy-Item .env.example .env
\end{lstlisting}

El archivo \texttt{.env} ser\'a utilizado por Docker Compose para proporcionar la
configuraci\'on necesaria a los tres servicios. Las variables previstas se describen
en la Tabla~\ref{tab:variables-entorno-construccion}.

\begin{table}[H]
\centering
\caption{Variables de entorno para la construcci\'on y ejecuci\'on del sistema.}
\label{tab:variables-entorno-construccion}
\begin{tabular}{|p{4.8cm}|p{8.2cm}|}
\hline
\textbf{Variable} & \textbf{Descripci\'on} \\
\hline
\texttt{POSTGRES\_DB} & Nombre de la base de datos que crear\'a PostgreSQL. \\
\hline
\texttt{POSTGRES\_USER} & Usuario administrador de PostgreSQL. \\
\hline
\texttt{POSTGRES\_PASSWORD} & Contrase\~na del usuario de PostgreSQL. \\
\hline
\texttt{DB\_HOST} & Nombre del servicio de base de datos dentro de Docker; su valor ser\'a \texttt{database}. \\
\hline
\texttt{DB\_PORT} & Puerto interno de PostgreSQL; por defecto, 5432. \\
\hline
\texttt{DB\_NAME} & Nombre de la base de datos utilizada por Express. \\
\hline
\texttt{DB\_USER} & Usuario empleado por el backend para acceder a PostgreSQL. \\
\hline
\texttt{DB\_PASSWORD} & Contrase\~na utilizada por el backend para acceder a PostgreSQL. \\
\hline
\texttt{PORT} & Puerto de escucha de la API Express; por defecto, 3000. \\
\hline
\texttt{GEMINI\_API\_KEY} & Clave de acceso al servicio Gemini utilizado por el asistente. \\
\hline
\texttt{MONTECARLO\_PYTHON\_PATH} & Ruta del ejecutable de Python dentro del contenedor del backend. \\
\hline
\texttt{EXPO\_PUBLIC\_API\_URL} & Direcci\'on de la API accesible desde el dispositivo m\'ovil. \\
\hline
\texttt{REACT\_NATIVE\_PACKAGER\_HOSTNAME} & Direcci\'on IPv4 de la m\'aquina anfitriona anunciada por Expo. \\
\hline
\end{tabular}
\end{table}

Las variables \texttt{POSTGRES\_DB}, \texttt{DB\_NAME},
\texttt{POSTGRES\_USER} y \texttt{DB\_USER} deber\'an mantener valores
coherentes. Lo mismo se aplica a \texttt{POSTGRES\_PASSWORD} y
\texttt{DB\_PASSWORD}, ya que el backend utilizar\'a esas credenciales para
conectarse a la base de datos creada por el contenedor.

Para utilizar Expo Go en un dispositivo f\'isico ser\'a necesario obtener la
direcci\'on IPv4 local de la m\'aquina anfitriona. En Windows puede consultarse
mediante \texttt{ipconfig}, mientras que en Linux o macOS puede utilizarse
\texttt{ip addr} o \texttt{ifconfig}. Si, por ejemplo, la direcci\'on obtenida es
\texttt{192.168.1.42}, la configuraci\'on incluir\'a los siguientes valores:

\begin{lstlisting}
DB_HOST=database
DB_PORT=5432
PORT=3000
MONTECARLO_PYTHON_PATH=python
EXPO_PUBLIC_API_URL=http://192.168.1.42:3000/api
REACT_NATIVE_PACKAGER_HOSTNAME=192.168.1.42
\end{lstlisting}

No se deber\'a utilizar \texttt{localhost} en
\texttt{EXPO\_PUBLIC\_API\_URL} cuando la aplicaci\'on se ejecute en un tel\'efono
f\'isico, ya que en ese caso \texttt{localhost} har\'ia referencia al propio
dispositivo y no al equipo que ejecuta Docker.

\subsection{Construcci\'on de las im\'agenes y arranque de los servicios}

Una vez configuradas las variables de entorno, la infraestructura completa se
construir\'a y ejecutar\'a desde el directorio ra\'iz mediante:

\begin{lstlisting}
docker compose up --build
\end{lstlisting}

Este comando construir\'a las im\'agenes del frontend y del backend, descargar\'a
la imagen de PostgreSQL, instalar\'a las dependencias npm y Python, crear\'a la red
interna y pondr\'a en marcha los tres contenedores. Al ejecutarse en primer plano,
la terminal mostrar\'a los registros de todos los servicios.

Para ejecutar la infraestructura en segundo plano se utilizar\'a la opci\'on
\texttt{-d}:

\begin{lstlisting}
docker compose up --build -d
\end{lstlisting}

El estado de los contenedores y sus registros podr\'an consultarse mediante:

\begin{lstlisting}
docker compose ps
docker compose logs -f backend
docker compose logs -f frontend
docker compose logs -f database
\end{lstlisting}

\subsection{Inicializaci\'on de PostgreSQL mediante el volcado}

La base de datos no requiere descargar conjuntos de datos externos ni ejecutar
scripts adicionales de extracci\'on. Toda la informaci\'on necesaria se encontrar\'a
en el volcado definitivo \texttt{database/init/dump.sql}.

El archivo \texttt{docker-compose.yml} montar\'a este volcado en el directorio
\texttt{/docker-entrypoint-initdb.d/} de la imagen oficial de PostgreSQL. De esta
forma, el esquema, las tablas y los datos se importar\'an autom\'aticamente la
primera vez que se cree el volumen de la base de datos. La importaci\'on inicial
puede comprobarse consultando los registros del servicio:

\begin{lstlisting}
docker compose logs -f database
\end{lstlisting}

Los archivos del directorio de inicializaci\'on solo se ejecutan cuando PostgreSQL
crea un directorio de datos vac\'io. Por tanto, si el contenedor ya dispone de un
volumen inicializado y se necesita cargar de nuevo el volcado, este podr\'a
importarse manualmente sin instalar PostgreSQL en la m\'aquina anfitriona:

\begin{lstlisting}
docker cp database/init/dump.sql database:/tmp/dump.sql
docker exec -it database psql -U <usuario> -d <base_de_datos> -f /tmp/dump.sql
\end{lstlisting}

En el segundo comando, \texttt{<usuario>} y \texttt{<base\_de\_datos>} deber\'an
sustituirse por los valores definidos en \texttt{POSTGRES\_USER} y
\texttt{POSTGRES\_DB}. Antes de repetir una importaci\'on sobre una base de datos
existente se deber\'a comprobar que no provoque duplicados o conflictos con los
objetos creados previamente.

\subsection{Puesta en marcha del backend}

El backend se iniciar\'a autom\'aticamente como parte de Docker Compose y quedar\'a
publicado en el puerto 3000 de la m\'aquina anfitriona. Express se conectar\'a al
servicio \texttt{database} a trav\'es de la red interna de Docker.

El entorno de Python tambi\'en formar\'a parte de este contenedor. Cuando el usuario
solicite una simulaci\'on manual, Express ejecutar\'a directamente el script
\texttt{simulacion\_montecarlo\_laliga.py} y recibir\'a su resultado en formato
JSON. No ser\'a necesario iniciar un servicio Python independiente.

Tras el arranque, la API y su documentaci\'on quedar\'an disponibles en las
direcciones indicadas en la Tabla~\ref{tab:servicios-expuestos}.

\begin{table}[H]
\centering
\caption{Servicios y direcciones del entorno local.}
\label{tab:servicios-expuestos}
\begin{tabular}{|p{3.4cm}|p{5.6cm}|p{4cm}|}
\hline
\textbf{Servicio} & \textbf{Direcci\'on} & \textbf{Funci\'on} \\
\hline
API REST & \texttt{http://localhost:3000/api} & Endpoints de la aplicaci\'on. \\
\hline
Documentaci\'on & \texttt{http://localhost:3000/api-docs} & Interfaz Swagger UI. \\
\hline
Metro Bundler & \texttt{http://<IP\_LOCAL>:8081} & Distribuci\'on de la aplicaci\'on Expo. \\
\hline
PostgreSQL & \texttt{database:5432} & Acceso exclusivo desde la red interna. \\
\hline
\end{tabular}
\end{table}

\subsection{Puesta en marcha del frontend}

El frontend tambi\'en se ejecutar\'a dentro de Docker, por lo que no ser\'a
necesario ejecutar \texttt{npm install} ni \texttt{npx expo start} directamente
en la m\'aquina anfitriona. El contenedor instalar\'a las dependencias y arrancar\'a
Expo con Metro Bundler en modo LAN y en el puerto 8081.

El c\'odigo QR generado por Expo podr\'a consultarse en los registros del
contenedor:

\begin{lstlisting}
docker compose logs -f frontend
\end{lstlisting}

Para abrir la aplicaci\'on se deber\'a iniciar Expo Go en el dispositivo m\'ovil y
escanear el c\'odigo QR. El tel\'efono y la m\'aquina anfitriona deber\'an estar
conectados a la misma red local, y el cortafuegos del sistema deber\'a permitir las
conexiones entrantes a los puertos 3000 y 8081.

Una vez cargado el c\'odigo JavaScript desde Metro, la aplicaci\'on se comunicar\'a
con Express mediante la direcci\'on definida en
\texttt{EXPO\_PUBLIC\_API\_URL}. El backend consultar\'a PostgreSQL y ejecutar\'a
el proceso Python cuando corresponda.

\subsection{Detenci\'on de la infraestructura}

Para detener los contenedores conservando el volumen y los datos de PostgreSQL se
utilizar\'a el siguiente comando:

\begin{lstlisting}
docker compose down
\end{lstlisting}

En el siguiente arranque, PostgreSQL reutilizar\'a el volumen existente y no
volver\'a a importar el volcado. La eliminaci\'on del volumen solo deber\'a
realizarse cuando se quiera reconstruir completamente la base de datos, puesto
que implica borrar de forma permanente los datos almacenados en el entorno local.
